\chapter{Monte Carlo Methods}

In the preceding chapter on machine learning, we saw that modern probabilistic
models often lead to integrals that cannot be evaluated with pencil and paper.

There is an old joke that the fastest way to compute a hard integral is to ask
a casino. Behind the joke is a practical principle: when an integral cannot be evaluated
in closed form, one can often \emph{simulate} it. By generating random samples,
the law of large numbers guarantees convergence to the true value. The result arrives not as
a formula but as an estimate -- a number accompanied by an error bar that shrinks as the sample
size increases.

These are the \emph{Monte Carlo methods}, named after the casino in
Monaco. The name was originally a wartime code word: the
entire subject is built on the idea that controlled randomness can be used as a computational tool. We will learn how to estimate quantities by simulation, how to
integrate functions that lack closed-form antiderivatives, how accurate those estimates are
via the central limit theorem, and how to improve efficiency
through variance reduction. The chapter concludes with Markov Chain Monte Carlo (MCMC),
a foundational method in modern Bayesian statistics and statistical physics.

Everything in this chapter rests on machinery you have already built. Monte
Carlo integration applies the law of large numbers directly to numerical calculation.
Its error bars are supplied by the central limit theorem. Its most advanced form, MCMC,
translates the theory of Markov chains into a computational algorithm.

\section{The Core Idea: Estimating by Simulation}
\label{ch35:core-idea}

Start with the simplest question there is. Suppose I want to know the
probability that a particular event happens.
Maybe it is a complicated event in a complicated system -- the probability that
a randomly shuffled deck has no card in its ``correct'' position, or that a
random network stays connected, or that a stock portfolio loses more than ten
percent in a year. The exact calculation might be a nightmare. But if I can
\emph{simulate} the system, I can just run it many times and count.

\begin{intuitionbox}
To estimate a probability, run the experiment $n$ times and report the fraction
of runs in which the event happened. To estimate an average, run the experiment
$n$ times and report the average of what you observed.
The mathematics is about how good these estimates are and how to make them
better.
\end{intuitionbox}

The simplest classic example is estimating $\pi$ by throwing darts. Draw the
unit square $[0,1]\times[0,1]$ and inscribe the quarter circle $x^2+y^2\le 1$.
Throw a dart uniformly at random into the square. The probability it lands
inside the quarter circle is the ratio of areas: $(\pi/4)/1 = \pi/4$. So if I
throw $n$ darts and a fraction $\hat p$ land inside, then $\hat p \approx \pi/4$,
which means $4\hat p \approx \pi$.

\begin{figure}[H]
\centering
\begin{tikzpicture}
\begin{axis}[
  width=8.2cm, height=8.2cm,
  xmin=0, xmax=1, ymin=0, ymax=1,
  axis equal image, axis lines=box,
  xtick={0,0.5,1}, ytick={0,0.5,1},
  title={$n=400$ darts; inside the arc estimates $\pi/4$},
  title style={font=\small}]
\addplot[densityred, very thick, domain=0:1, samples=120] {sqrt(1-x^2)};
\addplot[only marks, mark=*, mark size=0.6pt, sourceblue] table {
0.12 0.81 0.55 0.21 0.93 0.07 0.34 0.66 0.71 0.49
0.08 0.27 0.62 0.88 0.41 0.15 0.77 0.53 0.29 0.95
0.46 0.38 0.19 0.72 0.84 0.11 0.59 0.31 0.67 0.04
0.23 0.91 0.57 0.13 0.79 0.44 0.36 0.68 0.05 0.52
0.86 0.24 0.61 0.17 0.74 0.39 0.28 0.63 0.09 0.97
0.48 0.33 0.71 0.18 0.82 0.06 0.54 0.42 0.26 0.69
0.14 0.58 0.37 0.89 0.51 0.22 0.76 0.31 0.64 0.03
0.92 0.16 0.45 0.73 0.21 0.85 0.56 0.34 0.67 0.08
0.29 0.62 0.47 0.11 0.78 0.41 0.35 0.65 0.07 0.53
0.87 0.25 0.6 0.19 0.75 0.4 0.27 0.66 0.1 0.96
0.32 0.49 0.7 0.23 0.81 0.05 0.55 0.43 0.3 0.68
0.13 0.59 0.36 0.9 0.5 0.2 0.77 0.33 0.63 0.02
0.91 0.17 0.44 0.72 0.22 0.84 0.57 0.35 0.66 0.09
0.28 0.61 0.46 0.12 0.79 0.42 0.34 0.64 0.06 0.54
};
\end{axis}
\end{tikzpicture}
\caption{Estimating $\pi$ by Monte Carlo. Each dart is uniform on the square;
the fraction below the red arc estimates $\pi/4$. With $400$ darts you get one
or two correct digits -- accuracy improves only like $1/\sqrt{n}$.}
\label{ch35:fig-pi-darts}
\end{figure}
\begin{example}[Estimating $\pi$]
\label{ch35:ex-pi}
Suppose I throw $n = 10{,}000$ darts and $7{,}839$ land inside the arc. Then
$\hat p = 0.7839$ and my estimate is $\hat\pi = 4\hat p = 3.1356$. The true value
is $3.14159\dots$, so I am off by about $0.006$. Not bad -- but notice how much
work bought how little accuracy. Ten thousand darts gave two correct digits. To
get one more digit I would need roughly a hundred times as many darts. This slow
$1/\sqrt n$ convergence is the central fact, and the central frustration, of
Monte Carlo.
\end{example}
Why does this work at all? Because each dart is an independent trial, and the
indicator ``did this dart land inside?'' is a $\Bern(\pi/4)$ random variable. The
fraction inside is a sample mean of i.i.d.\ Bernoullis, and the law of large
numbers guarantees it converges to the true mean $\pi/4$.
\begin{historybox}
\textbf{Ulam, von Neumann, and Metropolis, 1946--1949: gambling on integrals.}
The modern Monte Carlo method was born at Los Alamos in the late 1940s. Stanislaw Ulam, recovering from an illness and playing games of Canfield solitaire, found the combinatorial equations for winning probabilities intractable. He realized he did not need a closed-form formula: he could simply \emph{deal out random hands and count the winners}.

He shared the insight with John von Neumann, who realized that new electronic digital computers (like ENIAC and MANIAC) could run such simulations for neutron-diffusion calculations in nuclear physics. Enrico Fermi had independently used similar manual sampling with a mechanical calculator in Rome years earlier. Nicholas Metropolis coined the code name ``Monte Carlo,'' after the Monaco casino where Ulam's uncle borrowed money to gamble. The method that began as a card game became one of the standard computational tools across physics, statistics, and machine learning.
\end{historybox}
\section{Monte Carlo Integration}
\label{ch35:integration}
The dart game is an example of numerical integration: the area under the arc is
$\int_0^1 \sqrt{1-x^2}\,\dd x = \pi/4$, which the random points estimated.
We now state the general principle, as high-dimensional integration is the
primary domain where Monte Carlo methods are essential.
Suppose we want to compute
\[
  I = \int_{\mathcal X} g(x)\, f(x)\, \dd x = \E[g(X)],
\]
where $X$ is a random variable (or vector) with density $f$, and $g$ is some
function. \emph{Any} expectation is an integral of this form, and conversely any
integral can be written as an expectation by choosing a convenient $f$. The
Monte Carlo recipe is about as simple as a recipe gets:
\begin{formalbox}
\textbf{Basic Monte Carlo estimator.} Draw $X_1,\dots,X_n$ independently from the
density $f$. Estimate
\[
  \hat I_n = \frac1n \sum_{i=1}^n g(X_i).
\]
By the law of large numbers, $\hat I_n \asto I$ as $n\to\infty$, provided
$\E\abs{g(X)} < \infty$.
\end{formalbox}
If you want a plain integral $\int_a^b h(x)\,\dd x$ with no probability in sight,
just rewrite it. Let $X\sim\Unif(a,b)$, whose density is $1/(b-a)$ on $[a,b]$.
Then
\[
  \int_a^b h(x)\,\dd x = (b-a)\int_a^b h(x)\,\frac{1}{b-a}\,\dd x
  = (b-a)\,\E[h(X)],
\]
so the estimator is $\hat I_n = \frac{b-a}{n}\sum_{i=1}^n h(X_i)$ with the
$X_i$ uniform on $[a,b]$. You sample points, evaluate the function, average, and
scale by the width. Done.
\begin{connectionbox}
Monte Carlo integration is a direct application of the law of large numbers.
Recall from Chapter~20 that for i.i.d.\ $Y_1,\dots,Y_n$ with mean $\mu$, the
sample mean $\bar Y_n \asto \mu$. Here $Y_i = g(X_i)$ and $\mu = \E[g(X)] = I$.
The reason Monte Carlo works in $1$ dimension or in
$1000$ dimensions, for smooth functions or rough ones, is that the LLN requires
no assumptions about dimension or smoothness, requiring only a finite mean.
This generality makes the method exceptionally versatile.
\end{connectionbox}
\subsection{The Error of Monte Carlo}
\label{ch35:error}
How accurate is $\hat I_n$? The central limit theorem supplies the convergence rate
directly, yielding one of the foundational formulas in scientific computing.
Let $\sigma^2 = \Var(g(X))$. The estimator $\hat I_n$ is a sample mean of i.i.d.\
terms, so it is unbiased, $\E[\hat I_n] = I$, and its variance is
\[
  \Var(\hat I_n) = \frac{\sigma^2}{n}.
\]
The standard error -- the typical size of the error -- is therefore
\[
  \SD(\hat I_n) = \frac{\sigma}{\sqrt n}.
\]
And by the CLT, for large $n$,
\[
  \hat I_n \approx \Normal\!\left(I, \frac{\sigma^2}{n}\right),
\]
so we can even put confidence intervals on our answer: roughly $95\%$ of the time
the truth lies within $\hat I_n \pm 1.96\,\sigma/\sqrt n$. In practice we do not
know $\sigma$, so we estimate it from the same sample,
\[
  \hat\sigma^2 = \frac{1}{n-1}\sum_{i=1}^n \bigl(g(X_i) - \hat I_n\bigr)^2,
\]
and report $\hat I_n \pm 1.96\,\hat\sigma/\sqrt n$. \emph{The simulation estimates
its own error bar.} Deterministic quadrature methods rarely provide an error estimate
so directly from function evaluations alone.
\begin{aestheticbox}
Notice what does \emph{not} appear in $\SD = \sigma/\sqrt n$: the
dimension $d$. A deterministic grid in $d$ dimensions with $m$ points per axis costs
$m^d$ evaluations -- the ``curse of dimensionality.'' Its error shrinks like
$m^{-2} = (m^d)^{-2/d}$, which becomes hopeless as $d$ grows. Monte Carlo's error
shrinks like $n^{-1/2}$ \emph{regardless of $d$}. In high dimensions this is not a
small advantage; it is the difference between possible and impossible. This is
why Monte Carlo, not quadrature, integrates the thousand-dimensional integrals of
Bayesian statistics and statistical physics.
\end{aestheticbox}

\begin{recoverybox}
\textbf{Pause and think: the $1/\sqrt{n}$ barrier.}
To gain a single additional decimal digit of precision in a Monte Carlo estimate, the standard error must shrink by a factor of $10$. By what factor must the number of simulated samples $n$ increase?
Because the error scales as $1/\sqrt{n}$, cutting the error by $10$ requires multiplying $n$ by $10^2 = 100$. Gaining three decimal digits requires a million times more computation! This quadratic penalty is why practitioners invest in clever variance reduction rather than relying on brute force alone.
\end{recoverybox}
\begin{figure}[H]
\centering
\begin{tikzpicture}
\begin{axis}[
  width=12cm, height=6.6cm,
  xlabel={number of samples $n$ (log scale)},
  ylabel={estimate of $\pi$},
  xmode=log, log basis x=10,
  xmin=10, xmax=100000, ymin=2.8, ymax=3.5,
  legend style={at={(0.92, 0.95)}, anchor=north east, font=\footnotesize, column sep=3pt},
  grid=both, grid style={gray!18}]
\addplot[densityred, very thick, domain=10:100000, samples=2] {3.14159265};
\addlegendentry{$\pi$}
\addplot[name path=hi, draw=none, domain=10:100000, samples=80]
  {3.14159265 + 1.96*1.642/sqrt(x)};
\addplot[name path=lo, draw=none, domain=10:100000, samples=80]
  {3.14159265 - 1.96*1.642/sqrt(x)};
\addplot[sourceblue, opacity=0.18, area legend] fill between[of=hi and lo];
\addlegendentry{$95\%$ band $\pm 1.96\sigma/\sqrt n$}
\addplot[structurepurple, thick, mark=none] coordinates {
(10,3.05)(20,3.27)(50,2.98)(100,3.21)(200,3.09)(500,3.165)
(1000,3.118)(2000,3.156)(5000,3.133)(10000,3.149)(20000,3.138)
(50000,3.1438)(100000,3.1408)};
\addlegendentry{running estimate}
\end{axis}
\end{tikzpicture}
\caption{A running Monte Carlo estimate of $\pi$ settling toward the truth. The
shaded $1/\sqrt n$ band narrows slowly: a hundredfold increase in
samples yields only a tenfold reduction in error.}
\label{ch35:fig-convergence}
\end{figure}
\begin{trapbox}
The error formula $\sigma/\sqrt n$ assumes $\sigma^2 = \Var(g(X))$ is finite. If
the integrand has heavy tails so that $\Var(g(X)) = \infty$, the CLT does not
apply, the estimator's fluctuations do not shrink like $1/\sqrt n$, and nominal
confidence intervals become invalid. Always verify whether the function being
averaged has a finite variance. A symptom in practice: the running estimate makes large
jumps that never settle, and a few outlier samples dominate the average.
\end{trapbox}
\begin{example}[A high-dimensional integral]
\label{ch35:ex-highdim}
Estimate $I = \int_{[0,1]^{10}} \cos\!\bigl(\tfrac{\pi}{2}\,\overline{x}\bigr)\,\dd x$,
where $\overline{x} = \frac{1}{10}\sum_{j=1}^{10} x_j$ is the average coordinate.
No grid can touch this: even a coarse $10$ points per axis is
$10^{10}$ evaluations. But Monte Carlo does not care about the $10$ dimensions.
Draw $n$ points uniformly in the cube, evaluate the cosine of the (scaled)
coordinate average at each, and average. With $n = 10^5$ samples you get three
reliable digits and an automatic error bar of size $\approx \hat\sigma/\sqrt{n}$,
all in a fraction of a second. The dimension never entered the error.
\end{example}
\section{Variance Reduction: Working Smarter}
\label{ch35:variance-reduction}
The error is $\sigma/\sqrt n$. We can shrink it two ways: increase $n$ (by drawing
more samples) or decrease $\sigma$ (by structuring the estimator more effectively). The goal of
\emph{variance reduction} is to find an estimator with the same mean $I$ but a
smaller variance. A method that halves $\sigma$ achieves the same precision as quadrupling $n$ -- without
increasing the sample size. We study two classic techniques: importance sampling
and antithetic variates.
\subsection{Importance Sampling}
\label{ch35:importance}
When sampling uniformly to estimate the expectation of $g$, one often evaluates
regions where $g$ is negligible and contributes little to the integral.
Efficiency improves by concentrating samples in regions where $g$ is large and
adjusting for the non-uniform sampling afterward. This adjustment is importance sampling.
\begin{intuitionbox}
Suppose you want to estimate the average wealth of a population. A simple uniform
sample rarely includes high-net-worth individuals, whose values significantly influence
the population mean. This leads to high estimator variance. Instead, you can
deliberately oversample higher wealth brackets and then down-weight their contributions
proportionally. Concentrating samples where the variation is greatest lowers overall variance.
That is the rationale behind importance sampling.
\end{intuitionbox}
Mathematically, suppose we want $I = \E_f[g(X)] = \int g(x) f(x)\,\dd x$ but we
sample from a different density $q$, called the \emph{proposal} or \emph{importance
distribution}. As long as $q(x) > 0$ wherever $g(x)f(x) \ne 0$, we can write
\[
  I = \int g(x)\,f(x)\,\dd x
    = \int \underbrace{\frac{g(x)\,f(x)}{q(x)}}_{\text{weighted integrand}} q(x)\,\dd x
    = \E_q\!\left[\frac{g(X)\,f(X)}{q(X)}\right].
\]
So the estimator becomes
\[
  \hat I_n^{\,\mathrm{IS}} = \frac1n \sum_{i=1}^n g(X_i)\,w(X_i),
  \qquad w(x) = \frac{f(x)}{q(x)},
  \qquad X_i \sim q.
\]
The weights $w(X_i)$ undo the fact that we drew from $q$ instead of $f$. The
estimator is still unbiased: $\E_q[g(X)w(X)] = I$. What changes is the variance.
\begin{whybox}
Why can importance sampling help? The variance of the new estimator is governed
by $\E_q[(gw)^2] = \int \frac{g(x)^2 f(x)^2}{q(x)}\,\dd x$. We get to choose $q$.
If we pick $q$ proportional to $\abs{g(x)}f(x)$ -- placing samples where
the integrand is largest -- this quantity decreases substantially. In the idealized
limit $q(x) \propto g(x) f(x)$ (assuming $g \ge 0$), the weighted integrand
$g(x)f(x)/q(x)$ becomes constant, yielding zero variance. A single
sample would give the exact answer. While this optimum cannot be achieved in practice
(since normalizing $q$ requires knowing $I$), it establishes the design principle:
\textbf{match the proposal density to the shape of the integrand.}
\end{whybox}
\begin{trapbox}
Importance sampling requires care. If $q$ is chosen poorly -- for example, with tails
lighter than $f$ -- an occasional sample can fall in a region where $q(x)$ is tiny relative
to $f(x)$, generating an enormous weight $f/q$. These rare extreme weights then dominate
the sample average, inflating the variance (sometimes to infinity) and making the estimator
far less reliable than standard Monte Carlo. A standard rule of thumb: the proposal $q$ must
have \emph{heavier tails} than the target $f$, never lighter.
\end{trapbox}
\begin{example}[Estimating a rare-event probability]
\label{ch35:ex-rare}
Suppose $X \sim \Normal(0,1)$ and we want $I = \P(X > 5)$. The true value is about
$2.87\times 10^{-7}$. Standard Monte Carlo draws standard normals and counts how many
exceed $5$; in ten million draws one expects only about three positive hits, so the
relative error is very large.
Importance sampling resolves this issue. Sample instead from a shifted normal,
$q = \Normal(5,1)$, which concentrates probability mass near $5$
(Figure~\ref{ch35:fig-importance}), so roughly half the
samples exceed $5$. Each sample carries weight $w(x) = \phi(x)/\phi(x-5)$, where
$\phi$ is the standard normal density. The estimator
$\frac1n\sum \ind\{X_i > 5\}\,w(X_i)$ converges with small relative error
using thousands of samples rather than billions. The sampler focuses on the
region of interest, while the weights maintain unbiasedness.
\end{example}
\begin{figure}[H]
\centering
\begin{tikzpicture}
\begin{axis}[
  width=12cm, height=5.6cm,
  xlabel={$x$}, ylabel={density},
  xmin=-3, xmax=8, ymin=0, ymax=0.45,
  legend pos=north east, samples=120, domain=-3:8]
\addplot[densityred, very thick] {exp(-x^2/2)/sqrt(2*pi)};
\addlegendentry{target $f=\Normal(0,1)$}
\addplot[sourceblue, thick, dashed] {exp(-(x-5)^2/2)/sqrt(2*pi)};
\addlegendentry{proposal $q=\Normal(5,1)$}
\addplot[name path=fcurve, draw=none, domain=5:8, samples=40]
  {exp(-x^2/2)/sqrt(2*pi)};
\addplot[name path=axisline, draw=none, domain=5:8, samples=2] {0};
\addplot[trapred, opacity=0.30] fill between[of=fcurve and axisline];
\draw[trapred, thick, ->] (axis cs:6.2,0.10) -- (axis cs:5.4,0.005);
\node[trapred, anchor=west, font=\small] at (axis cs:6.2,0.12)
  {rare region $\{x>5\}$};
\end{axis}
\end{tikzpicture}
\caption{Importance sampling for a rare event. Plain Monte Carlo draws from $f$
(red) and almost never lands in the rare region $\{x>5\}$, so the shaded tail
probability is estimated by counting a handful of hits. The proposal $q$ (blue,
dashed) is centred on the rare region, so about half its draws land there; the
weights $w=f/q$ correct for the shift and keep the estimator unbiased.}
\label{ch35:fig-importance}
\end{figure}
\subsection{Antithetic Variates}
\label{ch35:antithetic}
The second trick exploits \emph{negative correlation}. Recall that for two random
variables,
\[
  \Var\!\left(\frac{Y_1 + Y_2}{2}\right)
  = \frac14\bigl(\Var Y_1 + \Var Y_2 + 2\Cov(Y_1,Y_2)\bigr).
\]
If $Y_1$ and $Y_2$ are independent, the covariance is zero and we get the usual
variance reduction by averaging. But if we can arrange $\Cov(Y_1,Y_2) < 0$ -- so
that when one estimate is above the mean, the other tends to be below it -- the deviations
partially cancel, and the average has lower variance than under independence.
\begin{intuitionbox}
Consider combining two measurements whose errors are negatively correlated: whenever
one overestimates the quantity, the other tends to underestimate it. Averaging the two
yields smaller deviations overall. Antithetic variates implement this principle by pairing
each random draw with an inversely correlated counterpart.
\end{intuitionbox}
The standard construction: if you sample $U \sim \Unif(0,1)$, then $1-U$ is also
$\Unif(0,1)$, and the two are negatively correlated. Use the pair
$(U, 1-U)$ to generate two evaluations $g(U)$ and $g(1-U)$, average them, and
average those averages over many pairs. If $g$ is monotone increasing (or
decreasing), then $g(U)$ and $g(1-U)$ move in opposite directions, the covariance
is negative, and the variance drops significantly without requiring additional independent samples.
\begin{example}[Antithetic pairs in action]
\label{ch35:ex-antithetic}
Estimate $I = \int_0^1 e^x\,\dd x = e - 1 \approx 1.71828$. Plain Monte Carlo with
$U\sim\Unif(0,1)$ averages $e^U$; the variance of $e^U$ is
$\Var(e^U) = \E[e^{2U}] - (\E e^U)^2 = \frac{e^2-1}{2} - (e-1)^2 \approx 0.2420$.
Now use antithetic pairs: average $\frac12(e^U + e^{1-U})$. Because $e^x$ is
increasing, $e^U$ and $e^{1-U}$ are negatively correlated. A computation gives
$\Var\bigl(\tfrac12(e^U+e^{1-U})\bigr) \approx 0.0039$ -- about \emph{sixty times}
smaller than the plain estimator's per-sample variance. But because each antithetic
pair requires two function evaluations, the variance reduction for the same number
of function evaluations is about \emph{thirty times} (comparing $0.0039$ to $0.2420/2 \approx 0.1210$).
Thirty times less variance is the same accuracy as thirty times more brute-force samples,
for no extra cost.
\end{example}
\begin{happeningbox}
Both variance-reduction methods operate by allocating sampling effort strategically.
Standard Monte Carlo samples uniformly according to the base density. Importance sampling
concentrates draws where the integrand is largest in magnitude. Antithetic variates pairs draws
so that deviations negatively correlate and cancel upon averaging. Stratified sampling
partitions the domain to guarantee coverage across regions. In each case, one improves
efficiency by structuring the estimator rather than simply increasing the sample size.
\end{happeningbox}
\section{Generating Samples: Rejection Sampling}
\label{ch35:rejection}
The methods above assumed that draws from the reference distribution are readily
available. Standard algorithms exist for uniform and normal distributions. But many
distributions -- such as Bayesian posteriors defined only up to an unknown normalizing
constant -- cannot be sampled directly by standard generators. Rejection sampling
provides a general method for such settings, based on a direct geometric principle.
\begin{intuitionbox}
To sample a point uniformly from \emph{under the curve} of a density, throw darts
uniformly at a box that contains the curve, and keep only the darts that land
under the curve. The $x$-coordinates of the kept darts are distributed exactly
according to the density. Sampling from a density is equivalent to sampling
uniformly from the area beneath it.
\end{intuitionbox}
This rests on a clean fact worth stating on its own.
\begin{theorem}[Fundamental theorem of simulation]
\label{ch35:thm-fts}
Let $f$ be a density on $\R$. If $(X, Y)$ is uniformly distributed on the region
$\set{(x,y) : 0 \le y \le f(x)}$ beneath the graph of $f$, then the marginal
distribution of $X$ has density $f$.
\end{theorem}
\begin{proof}
The region under the graph has area $\int f(x)\,\dd x = 1$, so the uniform
density there is just $1$ on the region. The marginal density of $X$ is obtained
by integrating out $y$:
\[
  f_X(x) = \int_0^{f(x)} 1 \, \dd y = f(x). \qedhere
\]
\end{proof}
So if we can sample uniformly under the curve, we are done. Rejection sampling
does this when the curve is awkward by enclosing it in something easy.
\begin{formalbox}
\textbf{Rejection sampling.} Suppose we can evaluate a target density $f$ and we
have a \emph{proposal} density $q$ that we \emph{can} sample, together with a
constant $M$ with $f(x) \le M\,q(x)$ for all $x$ (the scaled proposal
``envelopes'' the target). Repeat:
\begin{enumerate}
  \item Draw $X \sim q$ and an independent $U \sim \Unif(0,1)$.
  \item If $U \le \dfrac{f(X)}{M\,q(X)}$, \textbf{accept} $X$; otherwise reject and
        try again.
\end{enumerate}
The accepted values are exact draws from $f$. The probability of accepting on a
given try is $1/M$, so on average it takes $M$ tries per accepted sample.
\end{formalbox}
\begin{figure}[H]
\centering
\begin{tikzpicture}
\begin{axis}[
  width=12cm, height=6.2cm,
  xlabel={$x$}, ylabel={density},
  xmin=-4, xmax=4, ymin=0, ymax=0.62,
  legend pos=north east, samples=120, domain=-4:4]
\addplot[densityred, very thick] {0.55*exp(-(x+1.2)^2/0.5)/sqrt(2*pi*0.25)*0.25
  + 0.45*exp(-(x-1.4)^2/0.8)/sqrt(2*pi*0.4)*0.4};
\addlegendentry{target $f(x)$}
\addplot[sourceblue, thick, dashed] {0.58*exp(-x^2/8)};
\addlegendentry{envelope $M\,q(x)$}
\addplot[only marks, mark=*, targetgreen, mark size=1.6pt] coordinates {(-1.2,0.05)};
\node[targetgreen, anchor=east] at (axis cs:-1.35,0.05) {\small accept};
\addplot[only marks, mark=*, trapred, mark size=1.6pt] coordinates {(1.0,0.30)};
\node[trapred, anchor=west] at (axis cs:1.1,0.30) {\small reject};
\end{axis}
\end{tikzpicture}
\caption{Rejection sampling. We throw points uniformly under the easy envelope
$M\,q$. Points landing under $f$ (green) are kept; points in the gap between $f$
and $M\,q$ (red) are thrown away. Kept points are exact samples from $f$.}
\label{ch35:fig-rejection}
\end{figure}
\begin{whybox}
Why are the accepted points distributed as $f$? Drawing $X\sim q$ and
$U\sim\Unif(0,1)$ and forming $(X,\,U\cdot M q(X))$ produces a point distributed
\emph{uniformly under the envelope} $Mq$. The acceptance test $U \le f(X)/(Mq(X))$
is exactly the condition that this point lies under $f$. So accepted points are
uniform under $f$, and by Theorem~\ref{ch35:thm-fts} their $x$-coordinates have
density $f$.
\end{whybox}
\begin{trapbox}
The cost of rejection sampling is $M$ proposals per accepted sample, and $M$ must
satisfy $f \le Mq$ \emph{everywhere}. If the proposal $q$ is a poor match for the
target -- especially in the tails -- then $M$ must be huge, and you reject almost
everything. In high dimensions this is fatal: the envelope's volume typically
exceeds the target's by a factor that grows exponentially with dimension, so the
acceptance probability $1/M$ becomes astronomically small. Rejection sampling
works well in one or two dimensions, but becomes impractical in moderate or high dimensions.
This limitation motivated the development of Markov chain Monte Carlo.
\end{trapbox}
\section{Markov Chain Monte Carlo}
\label{ch35:mcmc}
Consider now the problem of sampling from a complex, high-dimensional
distribution, such as a Bayesian posterior over many parameters or the
Boltzmann distribution of a physical lattice. Rejection sampling is thwarted
by the curse of dimensionality, and direct independent sampling is generally impossible.

Markov chain Monte Carlo (MCMC) addresses this by replacing independent sampling
with a dependent sequence. One constructs an ergodic Markov chain whose unique
stationary distribution is the desired target distribution $\pi$. As the chain runs,
the sequence of states it visits converges in distribution to $\pi$, and time averages
along the trajectory converge to expectations under $\pi$.
\begin{connectionbox}
This is the payoff of everything you learned about Markov chains
(Chapters~23--24). There you started with a transition matrix $P$ and found
the stationary distribution $\pi$ satisfying $\pi P = \pi$.

MCMC runs that logic \emph{backwards}: you are \emph{given} an
intractable distribution $\pi$ (such as a high-dimensional continuous posterior),
and you must \emph{design} a transition rule $P$ whose unique
stationary distribution is that $\pi$. The ergodic theorem for Markov chains then guarantees
that time averages along a single long trajectory converge to expectations under $\pi$:
\[
  \frac1n \sum_{t=1}^n g(X_t) \asto \E_\pi[g(X)].
\]
\end{connectionbox}
The key design tool is \emph{detailed balance} (reversibility), which you met
when studying stationary distributions. A chain with transition density
$P(x\to y)$ has $\pi$ as its stationary distribution if
\[
  \pi(x)\,P(x \to y) = \pi(y)\,P(y \to x) \qquad \text{for all } x,y.
\]
Read it as a flow-conservation law: in equilibrium, the probability current from
$x$ to $y$ exactly balances the current from $y$ to $x$. If detailed balance
holds, no net probability moves anywhere, so $\pi$ stays put -- it is stationary.
Every MCMC algorithm provides a systematic way to construct a transition rule that satisfies
detailed balance with respect to your chosen $\pi$.
\subsection{The Metropolis--Hastings Algorithm}
\label{ch35:metropolis}
The Metropolis--Hastings algorithm requires $\pi$ only \emph{up to a normalizing constant} --
which is exactly the situation in Bayesian inference, where the posterior is proportional to
(likelihood)$\times$(prior), but the normalizing integral is intractable. The algorithm
circumvents this integral entirely because it only ever evaluates the \emph{ratio}
$\pi(y)/\pi(x)$, in which the unknown normalizing constant cancels.
\begin{formalbox}
\textbf{Metropolis--Hastings.} Choose a proposal density $q(y \given x)$ from
which you can sample. Start at $x_0$. For $t = 0, 1, 2, \dots$:
\begin{enumerate}
  \item \textbf{Propose} a candidate $y \sim q(\cdot \given x_t)$.
  \item \textbf{Compute the acceptance ratio}
        \[
          \alpha = \min\!\left(1,\;
          \frac{\pi(y)\, q(x_t \given y)}{\pi(x_t)\, q(y \given x_t)}\right).
        \]
  \item \textbf{Accept or reject}: draw $U \sim \Unif(0,1)$. If $U \le \alpha$,
        set $x_{t+1} = y$ (move); otherwise set $x_{t+1} = x_t$ (stay).
\end{enumerate}
The resulting chain has stationary distribution $\pi$.
\end{formalbox}
\begin{intuitionbox}
Picture exploring an unknown surface whose height represents $\pi$, seeking high-density
regions while visiting states in proportion to their probability mass. From the
current state, the sampler proposes a candidate move $y$. If the candidate has higher
probability density ($\pi(y) > \pi(x)$), the step is always accepted. If the candidate
has lower density, it is still accepted with probability $\pi(y)/\pi(x)$. This balance
between accepting upward moves and occasionally accepting downward moves ensures the chain
traverses the full distribution rather than remaining trapped at a local mode
(Figure~\ref{ch35:fig-mcmc-path}).
\end{intuitionbox}
\begin{whybox}
Why does this leave $\pi$ stationary? We check detailed balance. For $x \ne y$
the transition density is $P(x\to y) = q(y\given x)\,\alpha(x,y)$. Suppose, say,
that $\pi(y)q(x\given y) \le \pi(x)q(y\given x)$, so that
$\alpha(x,y) = \frac{\pi(y)q(x\given y)}{\pi(x)q(y\given x)}$ and, in the reverse
direction, $\alpha(y,x) = 1$. Then
\[
  \pi(x)\,P(x\to y) = \pi(x)\,q(y\given x)\,
     \frac{\pi(y)q(x\given y)}{\pi(x)q(y\given x)}
   = \pi(y)\,q(x\given y),
\]
while
\[
  \pi(y)\,P(y\to x) = \pi(y)\,q(x\given y)\cdot 1 = \pi(y)\,q(x\given y).
\]
The two sides are equal, so detailed balance holds, and $\pi$ is stationary. The
$\min$ and the ratio are engineered precisely to make this equality hold. And
because only the ratio $\pi(y)/\pi(x)$ appears, any normalizing constant in $\pi$
cancels, allowing MCMC to handle unnormalized posteriors directly.
\end{whybox}
\begin{figure}[H]
\centering
\begin{tikzpicture}
\begin{axis}[
  name=landPlot,
  width=12cm, height=5.2cm,
  title={MCMC: Exploring an unknown probability surface $\pi(x)$},
  title style={font=\small},
  xlabel={state space $x$},
  ylabel={target altitude $\pi(x)$},
  domain=-3.5:3.5, samples=140,
  xmin=-3.5, xmax=3.5, ymin=0, ymax=1.1,
  axis lines=left,
  x label style={font=\footnotesize},
  y label style={font=\footnotesize},
  tick label style={font=\scriptsize},
]
% Bimodal density surface
\addplot[distblue, very thick, fill=distblue!12] {0.55*exp(-(x+1.5)^2/0.8) + 0.85*exp(-(x-1.2)^2/1.0)} \closedcycle;

% State x0 at -2.4 (altitude ~ 0.2)
\fill[black] (axis cs:-2.4, 0.2) circle (2.2pt);
\node[below, font=\scriptsize] at (axis cs:-2.4, 0.15) {$x_0$};

% Move 1: Uphill proposal to y1 = -1.6 (altitude ~ 0.54) -> Accepted!
\draw[->, very thick, targetgreen] (axis cs:-2.4, 0.2) to[out=40,in=160] (axis cs:-1.6, 0.54);
\fill[targetgreen!80!black] (axis cs:-1.6, 0.54) circle (2.2pt);
\node[above left, font=\scriptsize, targetgreen!80!black, align=center] at (axis cs:-1.8, 0.55) {Step 1: uphill\\(accepted)};
\node[below, font=\scriptsize] at (axis cs:-1.6, 0.48) {$x_1$};

% Move 2: Cross-valley proposal to y2 = 0.5 (altitude ~ 0.52) -> Accepted downhill/cross
\draw[->, very thick, quotientorange] (axis cs:-1.6, 0.54) to[out=15,in=150] (axis cs:0.5, 0.52);
\fill[quotientorange!80!black] (axis cs:0.5, 0.52) circle (2.2pt);
\node[above, font=\scriptsize, quotientorange!80!black] at (axis cs:-0.4, 0.65) {Step 2: downhill ratio (accepted)};
\node[font=\scriptsize] at (axis cs:0.62, 0.38) {$x_2 = x_3$};

% Move 3: Bad proposal into deep tail y3 = 3.2 (altitude ~ 0.01) -> Rejected!
\draw[->, thick, dashed, trapred] (axis cs:0.5, 0.52) to[out=15,in=140] (axis cs:3.0, 0.03);
\draw[trapred, thick] (axis cs:2.95,0.08) -- (axis cs:3.05,-0.02);
\draw[trapred, thick] (axis cs:2.95,-0.02) -- (axis cs:3.05,0.08);
\node[anchor=north east, font=\scriptsize, trapred, align=right] at (axis cs:3.4, 0.75) {Step 3: proposal into far tail\\(rejected, stays at $x_2$)};
\end{axis}

% Bottom lane: Timeline of samples
\begin{scope}[yshift=-1.2cm]
\draw[thick, ->] (-3.5,0) -- (3.5,0) node[right, font=\footnotesize] {$x$};
\draw[gray, dashed] (-1.5,-0.3) -- (-1.5,0.3) node[above, font=\tiny] {mode 1};
\draw[gray, dashed] (1.2,-0.3) -- (1.2,0.3) node[above, font=\tiny] {mode 2};

\fill[black] (-2.4,0) circle (2pt) node[below, font=\tiny] {$x_0$};
\fill[targetgreen!80!black] (-1.6,0) circle (2pt) node[below, font=\tiny] {$x_1$};
\fill[quotientorange!80!black] (0.5,0) circle (2pt) node[below, font=\tiny] {$x_2, x_3$};

\draw[->, targetgreen, thick] (-2.4,0.12) -- (-1.6,0.12);
\draw[->, quotientorange, thick] (-1.6,0.12) -- (0.5,0.12);
\node[font=\scriptsize, align=center, text width=9cm] at (0,-0.7) {Trajectory visits each region with frequency proportional to local mass $\pi(x)$};
\end{scope}
\end{tikzpicture}
\caption{Metropolis--Hastings exploring a target probability density. At each step, a candidate $y \sim q(y\given x_t)$ is proposed. Steps with $\pi(y) \ge \pi(x_t)$ are always accepted (green). Downhill steps ($\pi(y) < \pi(x_t)$) are accepted with probability $\alpha = \pi(y)/\pi(x_t)$ (orange), preventing the chain from remaining trapped at a single local peak. Proposals into low-probability regions are rejected (red), causing the chain to stay put ($x_{t+1}=x_t$), which correctly accumulates sampling weight in high-density regions.}
\label{ch35:fig-mcmc-path}
\end{figure}

\begin{example}[Sampling a standard normal with a random walk]
\label{ch35:ex-mh-normal}
Let the target be $\pi(x) \propto e^{-x^2/2}$ (a standard normal, but pretend we
only know it up to the constant). Use a symmetric random-walk proposal
$y = x_t + \eps$ with $\eps \sim \Normal(0, s^2)$. Symmetry means
$q(y\given x) = q(x\given y)$, so those terms cancel and the acceptance ratio
simplifies to the pure Metropolis form
\[
  \alpha = \min\!\left(1, \frac{\pi(y)}{\pi(x_t)}\right)
         = \min\!\left(1, e^{-(y^2 - x_t^2)/2}\right).
\]
Steps toward $0$ (where $\pi$ is large) are always accepted; steps into the tails
are accepted only sometimes. Run for ten thousand iterations and a histogram of
the visited states reproduces the bell curve. The step size $s$ matters: too
small and the chain inches along, exploring slowly; too large and almost every
bold proposal is rejected, so the chain sticks. A classic rule of thumb targets
an acceptance rate near $23\%$ for high-dimensional random-walk Metropolis.
\end{example}

\begin{figure}[H]
\centering
\begin{tikzpicture}
\begin{axis}[
  width=12cm, height=5.6cm,
  xlabel={iteration $t$}, ylabel={state $x_t$},
  xmin=0, xmax=120, ymin=-3.4, ymax=3.4,
  grid=both, grid style={gray!15}]
\addplot[structurepurple, thick] coordinates {
(0,2.8)(2,2.8)(4,2.1)(6,1.7)(8,1.7)(10,0.9)(12,1.3)(14,0.4)(16,0.4)
(18,-0.2)(20,0.5)(22,0.5)(24,-0.3)(26,-0.9)(28,-0.6)(30,-0.6)(32,0.1)
(34,0.8)(36,0.3)(38,0.3)(40,-0.4)(42,-1.1)(44,-0.7)(46,-0.7)(48,0.2)
(50,0.7)(52,1.4)(54,0.9)(56,0.9)(58,1.6)(60,0.8)(62,0.2)(64,-0.5)
(66,-0.5)(68,-1.3)(70,-0.6)(72,0.0)(74,0.6)(76,0.6)(78,1.1)(80,0.5)
(82,-0.2)(84,-0.8)(86,-0.8)(88,-1.5)(90,-0.9)(92,-0.3)(94,0.4)(96,0.4)
(98,1.0)(100,0.5)(102,-0.1)(104,-0.7)(106,-0.7)(108,0.0)(110,0.7)
(112,1.3)(114,0.8)(116,0.8)(118,0.2)(120,-0.4)};
\addplot[densityred, dashed, thick, domain=0:120, samples=2]{0};
\end{axis}
\end{tikzpicture}
\caption{A trace of a random-walk Metropolis chain targeting $\Normal(0,1)$. The
chain wanders around the high-probability center, occasionally venturing into the
tails. Flat segments are rejected proposals -- the chain stayed put. Time averages
along this single path estimate expectations under the normal.}
\label{ch35:fig-mh-trace}
\end{figure}

\subsection{Gibbs Sampling}
\label{ch35:gibbs}

When the target lives in many dimensions, proposing good joint moves is hard.
Gibbs sampling is a clever special case of Metropolis--Hastings that sidesteps
the problem: instead of moving all coordinates at once, update them \emph{one at
a time}, each from its \emph{conditional} distribution given all the others.

\begin{intuitionbox}
Suppose you and a partner are decorating a room and cannot agree on everything at
once. So you take turns: you fix the wall color given the current furniture, then
your partner picks furniture given the wall color, then you re-pick the wall
given the new furniture, and so on. Each step you solve a small one-variable
problem, holding everything else fixed. This back-and-forth converges
to a draw from the full \emph{joint} distribution.
\end{intuitionbox}

\begin{formalbox}
\textbf{Gibbs sampling.} To sample a joint distribution
$\pi(x_1, \dots, x_d)$, cycle through the coordinates. At each step, replace
coordinate $j$ by a fresh draw from its full conditional:
\[
  x_j^{\text{new}} \sim \pi\bigl(x_j \,\big|\, x_1,\dots,x_{j-1},
     x_{j+1},\dots,x_d\bigr).
\]
Repeat for $j = 1, 2, \dots, d$ and loop. The stationary distribution of the
resulting chain is exactly $\pi$.
\end{formalbox}

Gibbs sampling is Metropolis--Hastings with a proposal that is the exact
conditional distribution; one can check that its acceptance probability is always
$1$. Nothing is ever rejected -- every proposed move is taken. This makes Gibbs
attractive whenever the full conditionals have nice closed forms, which happens
constantly in hierarchical Bayesian models built from conjugate pieces.

\begin{example}[Gibbs for a bivariate normal]
\label{ch35:ex-gibbs}
Let $(X, Y)$ be bivariate normal with standard marginals and correlation $\rho$.
The full conditionals are themselves normal -- a fact you can derive from the
joint density:
\[
  X \given Y = y \;\sim\; \Normal(\rho y,\ 1 - \rho^2), \qquad
  Y \given X = x \;\sim\; \Normal(\rho x,\ 1 - \rho^2).
\]
Gibbs sampling alternates: draw $X$ given the current $Y$, then $Y$ given the new
$X$, and repeat. After a short burn-in the pairs $(X_t, Y_t)$ trace out the
elliptical correlation cloud of the true bivariate normal
(Figure~\ref{ch35:fig-gibbs-cloud}). We sampled a
two-dimensional distribution using only one-dimensional normal draws -- no matrix
square roots, no rejection.
\end{example}

\begin{figure}[H]
\centering
\begin{tikzpicture}
\begin{axis}[
  width=8.6cm, height=8.6cm,
  axis equal image, axis lines=box,
  xlabel={$x$}, ylabel={$y$},
  xmin=-3.2, xmax=3.2, ymin=-3.2, ymax=3.2,
  xtick={-3,0,3}, ytick={-3,0,3},
  title={Gibbs samples of a bivariate normal, $\rho=0.8$},
  title style={font=\small}]
% scatter cloud roughly along the rho=0.8 axis
\addplot[only marks, mark=*, mark size=0.7pt, structurepurple] coordinates {
(0.4,0.6)(-1.1,-0.9)(1.8,1.5)(0.2,-0.3)(-0.7,-1.0)(2.1,1.9)(-1.6,-1.3)
(0.9,1.2)(-0.3,0.1)(1.3,0.8)(-2.0,-1.7)(0.6,0.9)(-0.9,-0.6)(1.5,1.6)
(0.1,0.4)(-1.3,-1.5)(2.4,2.0)(-0.5,-0.8)(0.8,0.5)(-1.8,-1.4)(1.0,1.3)
(-0.2,-0.5)(1.7,1.4)(-1.0,-1.2)(0.5,0.2)(-2.2,-1.9)(1.2,1.0)(-0.6,-0.3)
(0.3,0.7)(-1.4,-1.1)(2.0,1.7)(-0.8,-1.0)(1.1,0.9)(-0.1,0.2)(1.6,1.8)
(-1.7,-1.6)(0.7,0.4)(-0.4,-0.7)(1.4,1.1)(-1.2,-0.9)(2.3,2.1)(-0.9,-1.1)
(0.0,0.3)(-1.5,-1.3)(1.9,1.5)(-0.7,-0.4)(0.6,0.8)(-2.1,-1.8)(1.3,1.2)
(-0.3,-0.6)(0.9,0.7)(-1.1,-1.4)(1.8,1.6)(-0.5,-0.2)(0.4,0.5)(-1.9,-1.5)
(1.0,1.1)(-0.8,-0.9)(2.2,1.9)(-0.2,0.0)(1.5,1.3)(-1.3,-1.0)(0.7,1.0)};
\end{axis}
\end{tikzpicture}
\caption{Output of Gibbs sampling for a bivariate normal with correlation
$\rho=0.8$. Each point is one $(X_t,Y_t)$ pair produced by alternately drawing one
coordinate from its conditional given the other. The cloud is the tilted
correlation ellipse of the target -- recovered using only one-dimensional normal
draws.}
\label{ch35:fig-gibbs-cloud}
\end{figure}

\begin{trapbox}
MCMC samples are \emph{correlated}, and this changes how you read the output.
Consecutive states are not independent, so the effective sample size is smaller
than the number of iterations -- sometimes far smaller. Two practical
consequences. First, \emph{burn-in}: the chain starts at an arbitrary point and
requires an initial segment to reach its stationary regime and reduce dependence
on the initial value; these early iterations are discarded.
Second, the naive error bar $\hat\sigma/\sqrt n$ is \emph{too optimistic} because
it assumes independence; the true variance of an MCMC average involves the
autocorrelation and is larger. And if the target has well-separated modes, a
random-walk chain can get stuck in one mode for a very long time, silently
missing the others -- a failure that no single-chain diagnostic reliably catches.
Run multiple chains from different starts and check that they agree.
\end{trapbox}

\begin{paradoxbox}
\textbf{The convergence you cannot see.} Here is the difficulty with MCMC.
The theory guarantees that time averages converge to the expectation under $\pi$
as $n \to \infty$, but the rate depends on the mixing time of the chain.
A chain trapped in one mode of a multimodal distribution can appear stationary:
the trace oscillates, the histogram looks smooth, and running averages settle, yet
the estimate misses other modes entirely. This distinguishes i.i.d.\ Monte Carlo,
whose standard error is governed directly by the central limit theorem, from MCMC,
whose finite-time accuracy depends on mixing throughout the state space.
The practical resolution relies on careful diagnostic procedure: run several
independent chains initialized from dispersed starting values, compare their output
(the Gelman--Rubin diagnostic formalizes this across-chain variance comparison),
and examine trace plots and autocorrelation times. MCMC guarantees asymptotic convergence,
but diagnosing finite-sample mixing requires systematic empirical verification.
\end{paradoxbox}

\section{Applications Across Science and Machine Learning}
\label{ch35:applications}

\begin{applicationbox}
\textbf{Bayesian inference.} The single most important use of MCMC. A Bayesian
posterior is
$\pi(\theta \given \text{data}) \propto p(\text{data}\given\theta)\,p(\theta)$,
and the constant of proportionality is an integral over all of parameter space
that is almost never computable. Metropolis--Hastings and Gibbs need only the
unnormalized product, so they let you draw from the posterior anyway. Tools like
Stan, PyMC, and JAGS are, underneath, automated MCMC engines. Every modern
hierarchical model -- in ecology, astronomy, pharmacology, sports analytics -- is
fit this way. When you read ``we used $4$ chains of $2000$ iterations,'' that is
MCMC.
\end{applicationbox}

\begin{applicationbox}
\textbf{Statistical physics.} Monte Carlo was born for physics and still lives
there. The Boltzmann distribution $\pi(s) \propto e^{-E(s)/kT}$ over the
configurations $s$ of a system (say, the spins of the Ising model of magnetism)
has an utterly intractable normalizing constant, the partition function. The
Metropolis algorithm -- proposed in that very 1953 paper -- samples
configurations to estimate magnetization, heat capacity, and phase transitions.
Simulating how a magnet behaves near its critical temperature is a Monte Carlo
calculation.
\end{applicationbox}

\begin{connectionbox}
\textbf{From interacting particles to Bayesian posteriors.}
The 1953 paper by Nicholas Metropolis, Arianna Rosenbluth, Marshall Rosenbluth,
Augusta Teller, and Edward Teller introduced what we now call the Metropolis algorithm
to compute equations of state for interacting particles in statistical mechanics.
In 1970, W.\ K.\ Hastings generalized it to asymmetric proposals, giving the
Metropolis--Hastings form we use today. For decades the method lived mostly in physics.
Then in the early 1990s, statisticians (notably Alan Gelfand and Adrian Smith, who
popularized Gibbs sampling for Bayesian models) demonstrated that it offered a practical
computational framework for Bayesian inference. Models that had been intractable
analytically could now be sampled and estimated numerically. A method originating in
statistical physics became a standard foundation for computational statistics and machine learning.
\end{connectionbox}

\begin{applicationbox}
\textbf{Machine learning.} Monte Carlo is everywhere in modern ML. Variational
inference and the ``reparameterization trick'' in variational autoencoders are
Monte Carlo estimates of intractable expectations over latent variables.
Policy-gradient methods in reinforcement learning estimate the gradient of
expected reward by sampling trajectories -- pure Monte Carlo, often with variance
reduction via ``baselines'' (a control-variate trick). Bayesian neural networks
quantify uncertainty with MCMC or its scalable cousins. Diffusion models, which
generate images, are built on a Markov chain that gradually denoises noise into
data. Dropout at test time can be read as a Monte Carlo estimate of predictive
uncertainty. The thread connecting all of these: an expectation no one can
compute, estimated by averaging samples.
\end{applicationbox}

\begin{applicationbox}
\textbf{Finance and risk.} Pricing a complex financial derivative means computing
$\E[\text{payoff}]$ under a model of how prices move -- usually a
high-dimensional, path-dependent expectation with no closed form. Banks estimate
it by simulating thousands of price paths and averaging the discounted payoffs,
with antithetic variates and importance sampling to tighten the estimate. Value
at Risk and stress testing simulate portfolio losses to estimate the probability
and size of rare, catastrophic outcomes -- exactly the rare-event regime where
importance sampling is effective (Example~\ref{ch35:ex-rare}).
\end{applicationbox}

\begin{aestheticbox}
A single principle -- replacing an intractable integral with a sample average
and relying on limit theorems for theoretical guarantees -- connects these diverse
applications. The law of large numbers guarantees convergence; the central limit
theorem characterizes the rate of error reduction; and Markov chain theory enables
sampling from complex, high-dimensional distributions.
\end{aestheticbox}

\section{Exercises}
\label{ch35:exercises}

\begin{exercise}[Warm-up mental checks]
\label{ch35:exr-warmup}
Answer these quick checks before reaching for paper:
\begin{enumerate}[label=(\alph*)]
  \item If you increase the number of Monte Carlo samples from $1{,}000$ to $100{,}000$, by what factor does the standard error shrink?
  \item In importance sampling, what happens to the weights and variance if your proposal distribution $q(x)$ has lighter tails than the target distribution $f(x)$?
  \item If you evaluate an integral using antithetic pairs $(U, 1-U)$, why does a strictly monotone integrand $g$ guarantee that the errors of the two evaluations tend to cancel?
  \item In the Metropolis algorithm, if the proposed candidate $y$ has higher target density than the current state $x_t$ ($\pi(y) > \pi(x_t)$), what is the probability of accepting the step?
\end{enumerate}
\end{exercise}

\begin{exercise}[Estimating $\pi$, and its error]
\label{ch35:exr-pi}
You throw $n = 4000$ darts at the unit square and $3{,}131$ land inside the
quarter circle. (a) Give the Monte Carlo estimate of $\pi$. (b) Using
$\hat p = 3131/4000$, estimate the standard error of $\hat\pi = 4\hat p$.
(\emph{Hint:} $\hat p$ is a sample mean of Bernoullis, so its variance is
$\hat p(1-\hat p)/n$, and $\hat\pi = 4\hat p$.) (c) Roughly how many darts would
you need to cut the standard error in half?
\end{exercise}

\begin{exercise}[Monte Carlo integration by hand]
\label{ch35:exr-integral}
Consider $I = \int_0^2 x^3 \, \dd x$ (whose true value is $4$). (a) Write $I$ as
$(b-a)\,\E[h(U)]$ for $U \sim \Unif(0,2)$ and identify $h$. (b) Given five uniform
draws on $[0,2]$ -- $0.4, 1.7, 0.9, 1.2, 1.9$ -- compute the Monte Carlo estimate
$\hat I_5$. (c) Is your estimate above or below the truth? Such small-sample
estimates are noisy; explain what the law of large numbers promises as $n$ grows.
\end{exercise}

\begin{exercise}[Why the error does not depend on dimension]
\label{ch35:exr-dimension}
A friend claims that integrating in $50$ dimensions must be vastly less accurate,
per sample, than integrating in $1$ dimension. Explain precisely why the Monte
Carlo standard error formula $\sigma/\sqrt n$ contains no dimension. Then explain
the catch: in what sense \emph{can} high dimension still hurt a Monte Carlo
estimate? (\emph{Hint:} think about what controls $\sigma$.)
\end{exercise}

\begin{exercise}[Antithetic variates]
\label{ch35:exr-antithetic}
You want $I = \int_0^1 \frac{1}{1+x}\,\dd x = \ln 2$. (a) Explain why
$g(x) = 1/(1+x)$ is monotone on $[0,1]$, so that antithetic pairing $(U, 1-U)$
should reduce variance. (b) With a single uniform draw $U = 0.3$, compute both the
plain one-sample estimate $g(U)$ and the antithetic estimate
$\frac12\bigl(g(U) + g(1-U)\bigr)$. (c) Compare each to the true value
$\ln 2 \approx 0.6931$ and comment on which is closer.
\end{exercise}

\begin{exercise}[Importance sampling weights]
\label{ch35:exr-importance}
You estimate $I = \E_f[g(X)]$ by sampling from a proposal $q$ and weighting by
$w(x) = f(x)/q(x)$. (a) Show the estimator $\frac1n\sum g(X_i)w(X_i)$ with
$X_i \sim q$ is unbiased for $I$. (b) Explain in words why a proposal $q$ with
\emph{lighter} tails than $f$ can produce occasional gigantic weights, and why
this inflates the variance. (c) What property should a good importance proposal
have relative to the integrand $g\cdot f$?
\end{exercise}

\begin{exercise}[Rejection sampling efficiency]
\label{ch35:exr-rejection}
You sample from a target $f$ using a proposal $q$ with envelope constant $M = 8$
(that is, $f(x) \le 8\,q(x)$ for all $x$, and $8$ is the smallest such constant).
(a) What is the probability of accepting on a given attempt? (b) On average, how
many proposals are needed per accepted sample? (c) If a tighter envelope with
$M = 1.5$ were available, by what factor would the method speed up? Relate this to
why rejection sampling fails in high dimensions.
\end{exercise}

\begin{exercise}[Metropolis acceptance ratio]
\label{ch35:exr-mh-ratio}
Consider Metropolis--Hastings targeting $\pi(x) \propto e^{-x^2/2}$ with a
symmetric proposal. The chain is currently at $x_t = 1.0$ and proposes
$y = -0.5$. (a) Since the proposal is symmetric, write the acceptance ratio
$\alpha$ in terms of $\pi$ only. (b) Compute $\alpha$ numerically. (c) Is the
move to $y$ accepted with certainty, and why does the direction of the move
(toward $0$) make this unsurprising?
\end{exercise}

\begin{exercise}[Detailed balance, conceptual]
\label{ch35:exr-detailed-balance}
State the detailed-balance equation for a Markov chain with stationary
distribution $\pi$ and transition rule $P(x\to y)$. (a) Explain in words why
detailed balance implies $\pi$ is stationary -- interpret it as a statement about
probability ``currents'' between states. (b) Is detailed balance necessary for
$\pi$ to be stationary, or only sufficient? (\emph{Hint:} recall that not every
stationary chain is reversible.)
\end{exercise}

\begin{exercise}[Gibbs for the bivariate normal]
\label{ch35:exr-gibbs}
Using the full conditionals from Example~\ref{ch35:ex-gibbs} with correlation
$\rho = 0.8$, suppose the chain is currently at $(X, Y) = (1.5, 0.0)$. (a) Write
the distribution from which you draw the new $Y$ given $X = 1.5$. (b) Give its
mean and variance numerically. (c) Explain why, as $\rho \to 1$, Gibbs sampling
mixes more and more slowly. (\emph{Hint:} how big are the conditional steps when
$1-\rho^2$ is tiny?)
\end{exercise}

\begin{exercise}[Diagnosing MCMC, proof-flavored]
\label{ch35:exr-diagnose}
Let $X_1, X_2, \dots$ be a stationary chain with $\Var(g(X_t)) = \sigma^2$ and
lag-$k$ autocorrelation $\rho_k = \Corr(g(X_t), g(X_{t+k}))$. (a) Show that
\[
  \Var\!\left(\frac1n\sum_{t=1}^n g(X_t)\right)
  = \frac{\sigma^2}{n}\left(1 + 2\sum_{k=1}^{n-1}\Bigl(1 - \tfrac{k}{n}\Bigr)\rho_k\right).
\]
(\emph{Hint:} expand the variance of the sum and collect terms by lag.) (b) When
all $\rho_k > 0$ (positive autocorrelation, typical of MCMC), is the true variance
larger or smaller than the i.i.d.\ value $\sigma^2/n$? (c) Explain why this means
the naive standard error $\hat\sigma/\sqrt n$ is overconfident for MCMC output,
and define what is meant by the ``effective sample size.''
\end{exercise}
